\section{Circuit RLC}

\subsection{Choix des composantes et explications}
\begin{frame}
	\frametitle{Choix des composantes}
	\begin{columns}
		\begin{column}{.5\textwidth}
			\begin{table}
				\caption{Composantes du RLC déterminées}
				\begin{tabular}{ll}
					\toprule
					Résistance & Valeur Choisie      \\
					\midrule\midrule
					$R_1$      & \SI{40}{\ohm}       \\
					$R_2$      & \SI{100}{\kilo\ohm} \\
					\bottomrule
				\end{tabular}
			\end{table}
		\end{column}
		\begin{column}{.5\textwidth}
			\begin{figure}
				\includegraphics[width=.7\textwidth]{figure/screenshot-rlc.jpg}
				\caption{Circuit du RLC}
				\label{fig:circuit-rlc}
			\end{figure}
		\end{column}
	\end{columns}
\end{frame}

\subsection{À la charge}
\begin{frame}
	\frametitle{Mises en équation de la charge}
	\begin{align}
		V_\mathrm{in}(t) & = V_C(t) + V_R(t) + V_L(t)
		\Rightarrow
		V_\mathrm{in}^{\prime\prime}(t) = V_L^{\prime\prime}(t) +
		\frac{R}{L}V_L^{\prime}(t) +
		\frac{1}{LC}V_L(t)                                                                    \\
		V_C(0)           & = 0;\quad
		V_R(0) = 0;\quad
		V_L(0) = 12;\quad
		I(0) = 0;                                                                             \\
		V_L(t)           & = \e^{-1000t} \left( 12\cos(8516.42t) - 1.41\sin(8516.42t) \right)
	\end{align}
\end{frame}

\begin{frame}
	\frametitle{Graphiques de charge théoriques}
	\begin{columns}

		\begin{column}{0.5\textwidth}
			\begin{figure}
				\includegraphics[width=.9\textwidth]{figure/vl_t_charge.jpg}
				\caption{M.E.É. du circuit RLC en charge}
				\label{graph:mee-charge-rlc}
			\end{figure}
		\end{column}

		\begin{column}{0.5\textwidth}
			\begin{figure}
				\includegraphics[width=.9\textwidth]{figure/tp12-dual.png}
				\caption{Simulation du circuit RLC en charge}
				\label{graph:lts-charge-rlc}
			\end{figure}
		\end{column}


	\end{columns}
\end{frame}

\subsection{À la décharge}
\begin{frame}
	\frametitle{Mises en équation de la décharge}
	\begin{align}
		V_\mathrm{in}(t) & = V_C(t) + V_R(t) + V_L(t)              \\
		V_\mathrm{in}(0) & = 0  ;\quad
		V_C(0) = 12;\quad
		V_R(0) = 0;\quad
		I(0) = 0;                                                  \\
		V_L^\prime(t)    & = -\frac{R}{L}V_L(t) = -\frac{R}{L}12 =
		60\times10^{6}                                             \\
		V_L(t)           & = -12 \e^{5\times10^{6}t}
	\end{align}
\end{frame}

\begin{frame}
	\frametitle{Graphiques de décharge théoriques}
	\begin{columns}

		\begin{column}{0.5\textwidth}
			\begin{figure}
				\includegraphics[width=.9\textwidth]{figure/vl_t_decharge.jpg}
				\caption{M.E.É. du circuit RLC en décharge}
				\label{graph:mee-decharge-rlc}
			\end{figure}
		\end{column}

		\begin{column}{0.5\textwidth}
			\begin{figure}
				
\includegraphics[width=.9\textwidth]{figure/todo.png}
				\caption{Simulation du circuit RLC en décharge}
				\label{graph:lts-decharge-rlc}
			\end{figure}
		\end{column}
	\end{columns}
\end{frame}

\subsection{Résultats expérimentaux}
\begin{frame}
	\frametitle{Résultats expérimentaux}
	\begin{columns}

		\begin{column}{0.5\textwidth}
			\begin{figure}
				
\includegraphics[width=.9\textwidth]{figure/todo.png}
				\caption{Résultats expérimentaux de la charge}
				\label{graph:osc-charge-rlc}
			\end{figure}
		\end{column}

		\begin{column}{0.5\textwidth}
			\begin{figure}
				
\includegraphics[width=.9\textwidth]{figure/todo.png}
				\caption{Résultats expérimentaux de la décharge}
				\label{graph:osc-decharge-rlc}
			\end{figure}
		\end{column}


	\end{columns}
\end{frame}
%vim: ft=tex
